\section*{Введение}
\addcontentsline{toc}{section}{Введение}

В современных условиях разработка программного обеспечения требует высокой скорости выпуска обновлений и постоянного контроля качества. Для решения этих задач команды используют практики непрерывной интеграции и непрерывной доставки, позволяющие автоматически запускать сборку, тестирование и развертывание приложений после изменений в исходном коде.

При росте количества проектов, окружений и команд CI/CD-процессы становятся сложнее. Информация о пайплайнах может находиться в разных системах: GitHub Actions, GitLab CI/CD, Jenkins, Argo CD и других инструментах. Разработчикам и DevOps-инженерам приходится переключаться между несколькими интерфейсами, вручную отслеживать статусы запусков, искать причины ошибок в логах и контролировать готовность релизов. Это увеличивает время реакции на сбои и снижает прозрачность процесса разработки.

Актуальность темы обусловлена необходимостью создания единого веб-интерфейса, который позволит централизованно управлять CI/CD-пайплайнами и отслеживать их состояние. Такое приложение может использоваться для просмотра проектов, запусков, этапов выполнения, логов, истории статусов и результатов сборки. Наличие единой панели мониторинга упрощает контроль качества и помогает быстрее выявлять проблемные участки процесса поставки программного обеспечения.

Цель курсового проекта --- разработать веб-приложение для управления и мониторинга CI/CD-пайплайнов.

Для достижения цели необходимо решить следующие задачи:
\begin{itemize}
    \item проанализировать предметную область и существующие CI/CD-инструменты;
    \item выбрать программные средства разработки и технологический стек;
    \item сформулировать требования к веб-приложению;
    \item спроектировать структуру приложения и модель данных;
    \item реализовать основные функции управления проектами, пайплайнами, запуском и просмотром логов;
    \item провести проверку работоспособности основных сценариев.
\end{itemize}

Объектом исследования являются процессы автоматизации сборки, тестирования и развертывания программного обеспечения. Предметом исследования является веб-приложение, обеспечивающее управление и мониторинг CI/CD-пайплайнов.

Курсовой проект состоит из введения, аналитического раздела, разделов проектирования и реализации, заключения и списка использованных источников. В первом разделе рассматриваются существующие CI/CD-системы, проводится выбор технологического стека и формулируются требования к разрабатываемому приложению.
